\chapter*{Sommario}
\addcontentsline{toc}{chapter}{Sommario}

Il mercato delle criptovalute si comporta come un sistema interconnesso: i rendimenti dei diversi asset tendono a muoversi insieme e la loro dipendenza si intensifica nei periodi di crisi. I modelli univariati ignorano per costruzione questa dipendenza, mentre il modello autoregressivo vettoriale (VAR) la rappresenta al prezzo di un numero di parametri che cresce quadraticamente con il numero di asset. La tesi si pone due questioni. La prima è quale vantaggio predittivo offra la struttura relazionale tra criptovalute, sfruttata da una Graph Neural Network, rispetto ad un approccio che le tratti come asset indipendenti; la seconda è come tale struttura si trasformi nel passaggio da un mercato stabile ad uno di crisi.

Il mercato è rappresentato come un grafo dinamico, ricostruito su una finestra mobile dalla correlazione dei rendimenti giornalieri di quindici criptovalute quotate su Binance tra il 2021 e il 2026. Una Graph Convolutional Network (GCN) è confrontata con modelli naive, autoregressivi e VAR e con una propria variante priva di grafo, tramite test di significatività statistica e un backtest al netto dei costi di transazione. L'evoluzione del grafo è misurata con metriche spettrali e topologiche attorno alla stretta regolamentare cinese del 2021, al crollo di Terra/Luna e al fallimento di FTX del 2022.

Quanto alla prima questione, nessun modello supera in modo significativo la previsione nulla e la GCN non si distingue dalla propria variante senza grafo. Il vantaggio che la strategia basata sulle sue previsioni mostra a costo nullo scompare già con una commissione di 10 punti base per cambio di posizione. Quanto alla seconda, il mercato risulta compatto per quasi tutto il periodo: oltre al modo comune a tutti gli asset, la matrice di correlazione non contiene struttura distinguibile dal rumore di stima. La stretta del 2021 porta quasi ogni metrica ad un estremo della propria distribuzione storica, mentre le crisi del 2022 le spostano solo marginalmente, perché investono un mercato già prossimo al suo stato di massima coesione. Nel campione analizzato la struttura relazionale coincide dunque in larga parte con il movimento comune del mercato e, alla frequenza giornaliera e ad un passo, non produce né un vantaggio predittivo né uno economico.
